\documentclass[10pt,a4paper]{amsart}
\usepackage[margin=19mm]{geometry}
\usepackage{amsmath,amssymb,mathtools}
\usepackage[scr=rsfs,cal=euler]{mathalfa}
\usepackage{xfrac}
\usepackage[unicode,hidelinks,bookmarks=false]{hyperref}
\newcommand{\ie}{i.e.\ }
\newcommand{\cf}{cf.\ }
\newcommand{\resp}{respectively\ }
\newcommand{\Subsection}{\subsection}
\providecommand{\nobreakdash}{\nobreak\hbox{-}}
\setlength{\emergencystretch}{2em}
\multlinegap0pt
\usepackage[normalem]{ulem}
\usepackage{xcolor}
\usepackage{amssymb}
\usepackage{enumerate}
\usepackage{tikz}
\newtheorem{theorem}{Theorem}
\title{Uniqueness of a topological Furstenberg system}
\author{Ioannis Kousek, Vicente Saavedra-Araya}
\date{Source: arXiv:2603.27899v1 (CC BY 4.0)}
\begin{document}
\maketitle

\noindent\emph{Source and licence.} Uniqueness of a topological Furstenberg system. Version arXiv:2603.27899v1 (CC BY 4.0), distributed by arXiv under the Creative Commons Attribution 4.0 International licence, http://creativecommons.org/licenses/by/4.0/, as recorded in the arXiv licence metadata for this exact version and shown on the abstract page https://arxiv.org/abs/2603.27899v1. Full licence notice: \texttt{licenses/arxiv-2603.27899-CC-BY-4.0.txt}.\par\smallskip

\noindent\textit{Editorial scope.} Counted unit: the article's theorem CHL (source0.tex lines 944--948). The article's complete original proof of it is delivered with it (source0.tex lines 949--985, one \texttt{proof} environment of 37 lines). No mathematics is written, repaired or re-derived: the statement and the proof are the article's own lines, in the article's own notation, and no retained line is substituted or re-flowed.  No further source-local context is needed: every cross-reference of every retained line resolves inside the two delivered spans.  Nothing was omitted from any retained span, so the package carries no omission record.  No retained line contains a citation, so the wrapper carries no bibliography and no bibliography entry.  No retained line contains a figure environment, an image-inclusion command or drawing code, so no image is delivered and the package declares no asset dependency.  Editorial formatting only, in addition: an AMS \texttt{amsart} preamble that restates the article's own declarations and macros used by the retained lines, namely the environments \texttt{theorem} on separate counters, together with the standard packages those lines need.  The retained environments print in this document as Theorem 1 in order of appearance, whereas the article prints the same items with the numbering of its own full document.  The complete native source of the article as distributed in the arXiv e-print for this exact version is bundled unchanged as \texttt{sources/197395/source0.tex}.
\begin{theorem}[Curtis-Hedlund-Lyndon]\label{CHL}
Let $G$ be a group, and let $f,f':G\to \mathcal{A}$ be finite-valued functions. Suppose $\phi$ is a factor map between the Furstenberg systems $(X_{R,f},(R_g)_g)$ and $(X_{R,f'},(R_g)_g)$. Then, there exists $g_1,\dots,g_k$ and a local code $\psi:\mathcal{A}^k\to \mathcal{A}$ such that
$$\phi(x)_h=\psi(x_{g_1h},\dots,x_{g_kh})$$
for all $x\in X_{R,f}$ and $h\in G$.
\end{theorem}

\begin{proof}
Let $a \in \mathcal{A}$, and suppose the cylinder set 
\[
C(a) = \{y \in X_{R,f'} : y_{e_G} = a\}
\]
is non-empty. By continuity of $\phi$, the preimage $\phi^{-1}(C(a))$ 
is open, and hence can be written as a (possibly uncountable) union of 
cylinder sets $(C'(a,i))_{i \in I_a}$ in $X_{R,f}$. 

Repeating this for each $a \in \mathcal{A}$ yields an open cover 
\[
(C'(a,i))_{a \in \mathcal{A},\, i \in I_a}
\]
of $X_{R,f}$. Since $X_{R,f}$ is compact, we can extract a finite 
subcover $(C_i)_{i \in I}$. Each $C_i$ depends only on finitely many 
coordinates and is contained in exactly one set of the form 
$\phi^{-1}(C(a))$. 

It follows that there exist $g_1, \dots, g_k \in G$ such that whenever 
$x_{g_i} = y_{g_i}$ for all $i = 1, \dots, k$, we have 
\[
\phi(x)_{e_G} = \phi(y)_{e_G}.
\]
Consequently, there exists a function $\psi : \mathcal{A}^k \to \mathcal{A}$ 
such that
\[
\phi(x)_{e_G} = \psi(x_{g_1}, \dots, x_{g_k}).
\]
Finally, since $R_h \circ \phi = \phi \circ R_h$ for all $h \in G$, we obtain that
\[
\phi(x)_h 
= (R_h \phi(x))_{e_G} 
= (\phi(R_h x))_{e_G} 
= \psi(x_{g_1 h}, \dots, x_{g_k h}).
\]

\end{proof} 

\end{document}
